% !TeX program = lualatex
% Standalone research notebook for original catalogue problem 61.
% Compile twice: lualatex 61.tex
% Original entry retained; research subsection and [E] references added.
% No external bibliography, image, data, or style file is required.
\documentclass[11pt,a4paper]{article}
\usepackage[margin=24mm,headheight=14pt]{geometry}
\usepackage{fontspec}
\setmainfont{Latin Modern Roman}
\setsansfont{Latin Modern Sans}
\setmonofont{Latin Modern Mono}
\usepackage{amsmath,amssymb,mathtools,amsthm}
\usepackage{unicode-math}
\setmathfont{Latin Modern Math}
\usepackage{microtype}
\usepackage{enumitem,xurl,needspace}
\usepackage[dvipsnames]{xcolor}
\usepackage{hyperref}
\hypersetup{unicode=true,colorlinks=true,linkcolor=MidnightBlue,
 urlcolor=MidnightBlue,bookmarksnumbered=true,
 pdftitle={Research notebook - Problem 61},
 pdfauthor={Research attempt; not independently refereed}}
\usepackage{bookmark}
\usepackage{titlesec}
\titleformat{\section}{\Large\bfseries\raggedright}{\thesection}{0.7em}{}
\usepackage{fancyhdr}
\pagestyle{fancy}\fancyhf{}
\fancyhead[L]{\small\sffamily Research notebook}
\fancyhead[R]{\small\sffamily Problem 61}
\fancyfoot[L]{\scriptsize 27 September 2026}
\fancyfoot[C]{\thepage}
\fancyfoot[R]{\scriptsize Unrefereed research attempt}
\setlength{\parindent}{0pt}
\setlength{\parskip}{4pt plus 1pt}
\setlength{\emergencystretch}{3em}
\clubpenalty=10000
\widowpenalty=10000
\displaywidowpenalty=10000
\setcounter{secnumdepth}{1}
\setlist[itemize]{leftmargin=3em,itemsep=2pt,topsep=3pt,parsep=0pt,before=\small}
\allowdisplaybreaks
\newcommand{\N}{\mathbb{N}}
\newcommand{\Z}{\mathbb{Z}}
\newcommand{\Q}{\mathbb{Q}}
\newcommand{\R}{\mathbb{R}}
\newcommand{\C}{\mathbb{C}}
\newcommand{\F}{\mathbb{F}}
\newcommand{\A}{\mathbb{A}}
\newcommand{\PP}{\mathbb P}
\newcommand{\E}{\mathbb{E}}
\newcommand{\eps}{\varepsilon}
\newcommand{\dd}{\,\mathrm d}
\newcommand{\sref}[2]{\hyperlink{source:#1:#2}{[S#2]}}
\newcommand{\eref}[2]{\hyperlink{extra:#1:#2}{[E#2]}}
\newif\ifshowcatalogue
\showcataloguetrue
\newcommand{\cataloguescope}[1]{\ifshowcatalogue
 \paragraph{Exact scope in the supplied catalogue.}
 \begin{quote}\small #1\end{quote}\fi}
\newtheorem{notebooklemma}{Lemma}[section]
\newtheorem{notebookproposition}[notebooklemma]{Proposition}
\numberwithin{equation}{section}
\begin{document}
\setcounter{section}{60}
% ================================================================
% problemId: problem.inverse-galois-problem-over-the-rational-numbers
% source releaseRank: 61; source releaseStatus: open_with_solved_subcases
% exactTarget SHA-256: 64fdd652fdff7ef1219176a515cdaa0512f55ed80c121d5378e9b9939fc2bb13
\clearpage
\section{\texorpdfstring{Inverse Galois problem over the rational numbers}{Inverse Galois problem over the rational numbers}}\label{61}
\subsection{Definitions and mathematical statement}
A finite extension $L/\Q$ is Galois when it is normal and separable. Its Galois group consists of field automorphisms of $L$ fixing $\Q$. The inverse problem is
\[
 \forall G\text{ finite}\ \exists L/\Q\text{ finite Galois}:\qquad
 \operatorname{Gal}(L/\Q)\simeq G.
\]
Only abstract group realization is demanded, not a uniform algorithm, a prescribed polynomial degree, or a regular realization over a rational-function field.
\par\smallskip\textit{Formulation and concept references:} \sref{61}{1}.
\cataloguescope{Prove or disprove that every finite group G is isomorphic to the Galois group Gal(L/Q) of some finite Galois extension L of the rational numbers Q.}
\subsection{Short English statement}
Does every finite group occur as the symmetry group of some finite Galois extension of the rationals?
% BEGIN RESEARCH INSERTION
\subsection{Research attempt}

\paragraph{Current frontier.}
Solvable-group constructions, Hilbert irreducibility, rigidity, and embedding problems provide major realization methods, with different hypotheses \sref{61}{1}, Sections II--III. The August 2026 manuscript of Huang--Jackson--Lee--Poonen--Pries--Zhang states an $M_{23}$ realization with an explicit degree-$23$ polynomial \eref{61}{1}, Theorem 1.1 and Example 1.2. This is relevant recent work, not a theorem for every finite group; its large exact computation is not independently repeated here. The deduction below does not depend on the new realization.

\paragraph{Attack route.}
A regular realization over $\Q(t)$ would suffice but is stronger than the target. Induction on a normal subgroup needs a proper solution of the corresponding embedding problem, not merely realization of its kernel and quotient. We instead combine a universal finite-group torsor with finitely many local Frobenius tests. The group-theoretic forcing will be proved; the global approximation step will be left visible.

\paragraph{Attempt: finite tests force the full group.}
Fix a nontrivial finite group $G$. Choose one representative $H$ from each conjugacy class of maximal proper subgroups.
\begin{notebooklemma}\label{r61:derangement}
There is $g_H\in G$ belonging to no conjugate of $H$. If $D\le G$ contains a conjugate of every selected $g_H$, then $D=G$.
\end{notebooklemma}
\begin{proof}
In the transitive action on $G/H$, double counting shows that the average number of fixed cosets is one. The identity fixes $[G:H]>1$ cosets, so some element fixes none. Such an element belongs to no conjugate of $H$. If $D$ were proper, it would lie in a conjugate of one of the chosen maximal subgroups, contradicting the prescribed element for that subgroup.
\end{proof}

Let $W=\A^{|G|}_{\Q}$ be the regular permutation representation and remove the fixed loci of all nonidentity elements. The free open locus $U$ is nonempty: a vector with distinct coordinates has trivial stabilizer. Restricting the finite affine quotient gives a finite étale torsor
\[
 q:U\longrightarrow V=U/G.
\]
A rational point $v\in V(\Q)$ defines, up to conjugacy, a continuous homomorphism $\rho_v:G_{\Q}\to G$. Its fiber is connected precisely when $\rho_v$ is onto; a connected fiber is then a Galois field extension with group $G$. These are the standard torsor and descent conventions \eref{61}{2}, Sections 4.4--4.5 and 5.12.

\paragraph{Local construction without a global realization.}
Choose distinct primes $p_H$. The unramified quotient $G_{\Q_{p_H}}\twoheadrightarrow\widehat{\Z}$ has a homomorphism to $G$ sending Frobenius to $g_H$, defining a local $G$-torsor $T_H$. Its associated étale algebra may be disconnected; that is allowed locally.

Every such $T_H$ occurs as a fiber of $q$. Twist the linear representation $W$ by $T_H$. Linear Galois descent identifies the twist with a vector space over $\Q_{p_H}$. Its twisted free locus is a nonempty Zariski-open subset and has a rational point, since an infinite field's affine space cannot be covered by finitely many proper polynomial zero sets. Untwisting gives a $G$-equivariant map $T_H\to U$ and hence a point $v_H\in V(\Q_{p_H})$ with the prescribed torsor fiber.

There is an analytic neighborhood $\mathcal O_H$ of $v_H$ on which the same $G$-torsor type persists. Krasner's lemma gives local constancy of the finite étale fiber; the local identifications can be chosen equivariantly by retaining the finitely many $G$-action maps. Alternatively the twist of $q$ is an étale map with a rational point, so its image contains an analytic neighborhood and its points trivialize the twisted torsor. This uses the local openness result in \eref{61}{2}, Proposition 3.5.73, as well as Proposition 3.5.74.

\begin{notebookproposition}\label{r61:approx}
For this quotient and these neighborhoods, suppose there exists $v\in V(\Q)$ with
\begin{equation}\label{r61:wa}
 v_{p_H}\in\mathcal O_H\qquad\text{for every chosen }H.
\end{equation}
Then $G$ occurs as a Galois group over $\Q$. Establishing this finite approximation assertion for every finite $G$ would resolve the target.
\end{notebookproposition}
\begin{proof}
The image of $\rho_v$ contains a conjugate of every $g_H$, by restriction to the selected decomposition groups. Lemma~\ref{r61:derangement} makes it all of $G$. The fiber is the desired field. The trivial group is separately realized by $\Q$.
\end{proof}

This asks for one specially chosen finite system of local conditions, not weak approximation for every neighborhood on every quotient. It remains a substantial arithmetic assertion, stronger than simply exhibiting a rational point of $V$.

\paragraph{Exact small-instance check.}
For $S_3$, take $f(t)=t^3-t-1$. It is irreducible modulo $2$, being a cubic without a root. Modulo $5$,
\[
 f(t)=(t-2)(t^2+2t+3),
\]
and the quadratic has nonsquare discriminant $2$. The integer discriminant is $-23$, so the reductions are unramified. The splitting-field group contains a $3$-cycle and a transposition and is therefore $S_3$. The script verifies these elementary modular calculations. This checks the forcing mechanism, not the general approximation assertion.

\paragraph{Stress test: a dominant quotient map does not patch torsors.}
Descending rational points of $U$ gives only trivial torsor fibers. Already for $G=C_2$, $U=\mathbb G_m$ and $q(x)=x^2$, these parameters are rational squares. They cannot enter a $p$-adic neighborhood of nonsquares. The quotient itself is rational in this example and has other rational points; the error is inferring the desired approximation from rational lifts.

The local points constructed at different primes have not been patched. Invoking Hilbert irreducibility would require a suitable rational parameter space with the needed specialization behavior, which has not been provided. Nor does the argument assume closure of realizable groups under arbitrary extensions, or infer arbitrary-group realization from realization of simple groups.

\paragraph{Outcome.}
\textbf{Outcome: Conditional reduction.}

We obtain a finite maximal-subgroup obstruction set, construct its local Frobenius tests on an explicit versal quotient, and prove that one rational point meeting the tests realizes $G$. The missing step is~\eqref{r61:wa} for every finite group. No general patching theorem or nonrealizable group is produced. The torsor ingredients are classical; the notebook records their assembly and exact remaining condition, without a novelty claim for the method.

% END RESEARCH INSERTION
\subsection{Sources}
\begin{itemize}\raggedright
\item[\textnormal{[S1]}]\hypertarget{source:61:1}{} Fariba Ranjbar and Saeed Ranjbar, arXiv:1512.08708 (2015).
\url{https://arxiv.org/abs/1512.08708}.
{\small\textit{Catalogue formulation source.}}
\item[\textnormal{[S2]}]\hypertarget{source:61:2}{} The Mathieu group M\_23 is a Galois group over Q.
\url{https://arxiv.org/html/2608.08538v1}.
{\small\textit{Catalogue research/status reference; not a proof endorsement.}}
\item[\textnormal{[S3]}]\hypertarget{source:61:3}{} The inverse Galois challenge, part II.
\url{https://galoisrepresentations.org/2026/08/16/the-inverse-galois-challenge-part-ii/}.
{\small\textit{Catalogue research/status reference; not a proof endorsement.}}

% BEGIN ADDED RESEARCH REFERENCES
\item[\textnormal{[E1]}]\hypertarget{extra:61:1}{} Xiaoyu Huang, Blake Jackson, Kyu-Hwan Lee, Bjorn Poonen, Rachel Pries, and Shaowu Zhang, \emph{The Mathieu group $M_{23}$ is a Galois group over $\Q$}, arXiv:2608.08538 (9 August 2026), Theorem 1.1 and Example 1.2.
\url{https://arxiv.org/abs/2608.08538}.
{\small\textit{Recent manuscript; its large explicit verification is not independently reproduced or used in the conditional proof. Consulted 27 September 2026.}}
\item[\textnormal{[E2]}]\hypertarget{extra:61:2}{} Bjorn Poonen, \emph{Rational Points on Varieties}, Graduate Studies in Mathematics 186, American Mathematical Society, 2017; Sections 4.4--4.5 and 5.12, Propositions 3.5.73--3.5.74.
\url{https://math.mit.edu/~poonen/papers/Qpoints.pdf}.
{\small\textit{Descent, twists, torsors, and local étale openness/Krasner lemma. Consulted 27 September 2026.}}
% END ADDED RESEARCH REFERENCES
\end{itemize}

\end{document}
